\section{Standard-type theta points in a declared additive carrier}
\label{sec:uniform-theta-points}

We now restrict the local source to numerical representatives of
standard-type theta values, with their finite root-of-unity ambiguity.
For each independent choice of these representatives, one actual
full-Galois Kummer operator attains every point-hull bound below.
The operator acts on a declared native additive carrier. This does
not identify its action there with an action on the original
multiplicative theta Kummer class, or establish membership in the
global output family of \cite[Corollary~3.12]{IUT3}.

Page references to the theta and IUT sources in this section use
the archived author versions: December 2008 for \cite{MochEtaleTheta},
May 2020 for \cite{IUT1,IUT3}, and December 2020 for \cite{IUT2}.

\subsection{The normalized evaluation and its finite ambiguity}

Retain the field and roots of \eqref{eq:gs-field-data}. Thus
$\ell\ge7$ and $p$ are primes, $p\equiv-1\pmod{30\ell}$,
$v_p(b_0)=2$, and
\begin{equation}\label{eq:utp-field}
 \begin{gathered}
 e=15\ell,\quad d=2e,\quad h=(\ell-1)/2,\quad
 K_0=\Q_p(\mu_{30\ell}),\\
 \varpi^e=b_0,\quad E=K_0(\varpi),\quad q=b_0^2,\quad
 r_0=\varpi^{15}=q^{1/(2\ell)},\\
 \beta=\varpi^{(e+1)/2}/p,\quad
 \kappa=1-1/e,\quad I=\beta^{1-e}\OO_E .
 \end{gathered}
\end{equation}
Here $\varpi$ is the element denoted $\pi$ in
\eqref{eq:gs-field-data}, and is not a uniformizer.
Lemma~\ref{lem:gs-field} gives
$I=p^{-1}\log\OO_E^\times$, $\operatorname{Tr}I=\Z_p$,
and $\mu_{2\ell}\subset K_0$. In particular, $e\le p-2$,
$p\nmid d$, and $E/\Q_p$ is Galois with residue degree two.

Choose $i^2=-1$ in an algebraic closure. The series of
\cite[Proposition~1.4, pp.~20--21]{MochEtaleTheta} is
\begin{equation}\label{eq:utp-series}
 \vartheta(U)=
 \sum_{n\in\Z}(-1)^nq^{n(n+1)/2}U^{2n+1},
 \qquad f(U)=\vartheta(U)/\vartheta(i).
\end{equation}
The standard-type theta function is the reciprocal of an
$\ell$-th root of this normalized function
\cite[Example~3.2(ii),(iv)]{IUT1}.

\begin{lemma}[The normalized theta coset]\label{lem:utp-evaluation}
For every integer $a$ and every nonzero $U$ in a finite extension
of $E$,
\[
 \vartheta(q^{a/2}U)
   =(-1)^a q^{-a^2/2}U^{-2a}\vartheta(U),
 \qquad
 f(\pm i q^{a/2})=\pm q^{-a^2/2}.
\]
The reciprocals of the $\ell$-th roots of these two evaluations
form exactly the coset
\begin{equation}\label{eq:utp-coset}
                         \mu_{2\ell}r_0^{a^2}.
\end{equation}
\end{lemma}
\begin{proof}
At the indicated points the term valuations grow quadratically
as $|n|\to\infty$, so reindexing is legitimate. Substitution
$n\mapsto n+a$ gives the functional equation. Also
$\vartheta(-U)=-\vartheta(U)$. At $U=i$ the terms $n=0,-1$
sum to $2i$, and all other terms have positive valuation.
Thus $\vartheta(i)$ is a unit. Since $i^{-2a}=(-1)^a$,
the two normalized evaluations follow. Their reciprocal roots
have $\ell$-th powers $\pm q^{a^2/2}$, whose complete root set,
for odd $\ell$, is \eqref{eq:utp-coset}.
\end{proof}

In the standard cyclotomic and Kummer marking,
\cite[Remark~2.5.1(i), p.~72]{IUT2} identifies the unperfected
evaluation set at label $j$ with the Kummer image of
$\mu_{2\ell}r_0^{j^2}$, where the exponent uses the integer
representative $0\le j\le h$. The ambiguities at different
labels are independent \cite[Remark~2.5.1(iii), p.~73]{IUT2}.
They are not arbitrary elements of $\OO_E^\times$:
standard-type normalization removes that larger constant-unit
ambiguity \cite[Remark~1.12.2(i),(ii)]{IUT2}.

\subsection{One arrow for independent torsion representatives}

For each $1\le j\le h$, choose $\zeta_j\in\mu_{2\ell}$
independently and put
\begin{equation}\label{eq:utp-inputs}
 \begin{gathered}
 x_j=\zeta_j r_0^{j^2},\qquad r_j=v_p(x_j)=2j^2/\ell,\\
 n_j=\lfloor r_j\rfloor,\qquad k_j=n_j+1,\qquad
 T_j=p^{-k_j}x_j,\qquad
 V=I/pI,\quad V_0=\beta I/pI .
 \end{gathered}
\end{equation}
The quotient $V/V_0$ has dimension two over $\mathbf F_p$.
Bars will denote reduction modulo $pI$.

\begin{lemma}[Content, trace and projective inertia]
\label{lem:utp-content}
For every choice of the $\zeta_j$,
\begin{equation}\label{eq:utp-content}
 \begin{gathered}
 k_j=\lfloor r_j+\kappa\rfloor,\qquad
 T_j\in\beta I\setminus pI,\qquad
 \operatorname{Tr}_{E/\Q_p}(T_j)=0,\\
 1\in\beta I\setminus pI,\qquad
 \operatorname{Tr}_{E/\Q_p}(1)=d\in\Z_p^\times .
 \end{gathered}
\end{equation}
If $\sigma(\beta)=\epsilon\beta$ generates
$\operatorname{Gal}(E/K_0)$, with $\epsilon$ of order $e$, then
\begin{equation}\label{eq:utp-character}
                         \sigma(T_j)=\epsilon^{30j^2}T_j .
\end{equation}
The projective orbit of $\overline{T_j}$ over $\mathbf F_p$
has exactly $\ell$ elements.
\end{lemma}
\begin{proof}
Write $2j^2=\ell n_j+\delta_j$, where
$1\le\delta_j\le\ell-1$. Then
$0<\delta_j/\ell-1/e<1$, proving the formula for $k_j$.
Moreover
\[
 e\,v_p(T_j)=15\delta_j-e,\qquad
 2-e<15-e\le15\delta_j-e\le-15<1 .
\]
The exponents $2-e$ and $1$ are the thresholds for $\beta I$
and $pI$. This proves the two membership assertions, including
the analogous assertion for $1$, whose exponent is zero.

With $\gamma=p^2/b_0\in\Z_p^\times$,
\eqref{eq:gs-uniformizer} gives $r_0=\gamma^{15}\beta^{30}$.
Both $\gamma$ and $\zeta_j$ lie in $K_0$ and are fixed by
$\sigma$, proving \eqref{eq:utp-character}. Its character
has order
$e/\gcd(e,30j^2)=\ell$. Summing its inertia conjugates gives
zero trace over $K_0$, and then zero absolute trace.

These roots reduce injectively and
$\mu_\ell\cap\mathbf F_p^\times=\{1\}$, since
$p\equiv-1\pmod\ell$. If two inertia images of
$\overline{T_j}$ are proportional over $\mathbf F_p$,
their character residues are therefore equal. Indeed, a
nonzero difference of residues is a unit of $K_0$, and
multiplying $T_j\notin pI$ by such a unit cannot put it in
$pI$. Thus proportionality forces the inertia exponents
to agree modulo $\ell$, while agreement gives equality
already in $E$. This proves the projective orbit assertion.
\end{proof}

Let $\Gamma_E$ be the actual integral Kummer operators of
Section~\ref{sec:full-galois-lifts}, extended $\Q_p$-linearly
to $E$. They preserve $I$; we do not identify this group of
operators with every integral linear automorphism.

\begin{theorem}[Common attainment for the finite theta source]
\label{thm:utp-common-arrow}
For every tuple $(\zeta_1,\ldots,\zeta_h)\in\mu_{2\ell}^h$,
there is one $F\in\Gamma_E$ such that
\begin{equation}\label{eq:utp-common}
 v_p(F(1))=v_p(F(T_j))=-\kappa
                      \qquad(1\le j\le h).
\end{equation}
The chosen $F$ may depend on the tuple.
\end{theorem}
\begin{proof}
The hypotheses of Lemma~\ref{lem:gl-integral-basis} hold:
$d$ is even, $2\le e\le p-2$, and $\operatorname{Tr}I=\Z_p$.
Use its basis
\[
 I=\Z_p a_{\rm JW}\oplus\Z_p b_{\rm JW}\oplus W,\qquad
 I\cap\ker\operatorname{Tr}=\Z_p a_{\rm JW}\oplus W,
\]
and let $B$ be the $b_{\rm JW}$-coordinate. The trace of
$b_{\rm JW}$ is a unit, so $B=\operatorname{Tr}/
\operatorname{Tr}(b_{\rm JW})$.

For each $j$, at most one of the first $\ell$ inertia powers
can place $\overline{T_j}$ on
$\mathbf F_p\overline a_{\rm JW}$. Since $h<\ell$,
one power avoids this line for every label. It fixes $1$
and preserves $V_0$ and trace. We now verify every remaining
hypothesis of Lemma~\ref{lem:gl-finite-family}, with
$k=\mathbf F_p$, $L=V_0$, $U=\overline1$ and the transformed
vectors $\overline{T_j}$. All lie in $L$ by
Lemma~\ref{lem:utp-content}; $B(U)\ne0$, $B(T_j)=0$,
and $T_j\notin\mathbf F_p a_{\rm JW}$.
The trace kernel maps onto $V/V_0$: for $y\in I$,
subtract $\operatorname{Tr}(y)/d\in\OO_E\subset\beta I$.
Finally its cardinality condition is $h+1<p$.

That lemma consequently gives one parabolic composition
sending all these vectors outside $V_0$. Its integer-coordinate
crosses and symplectic operations have actual full-group lifts
by Proposition~\ref{prop:gl-full-descent} and
Lemma~\ref{lem:gl-central}, as used in the arithmetic conclusion
of Lemma~\ref{lem:gl-finite-family}. Compose them with the
chosen field inertia automorphism. The resulting canonical
action $M$ sends every input into $I\setminus\beta I$, of
valuation $-\kappa$.

Proposition~\ref{prop:gl-variance} gives Kummer action
$c_\alpha M_\alpha^{-1}$, with $c_\alpha\in\Z_p^\times$.
Using the inverse of the constructed group automorphism
therefore gives a unit multiple of $M$, proving
\eqref{eq:utp-common}.
\end{proof}

This proves $\forall(\zeta_j)\,\exists F\,\forall j$,
not the stronger assertion that one $F$ works for every
tuple. It does not commute multiplication by $\zeta_j$
with $F$, which need not be $K_0$-linear.

\subsection{The additive point hull and the logarithmic bridge}

For $m=m_j=j+1$, let
\[
 \mathcal T_m=E^{\otimes_{\Q_p}m},\qquad
 A_m=\OO_E^{\otimes_{\Z_p}m},\qquad
 B_m=\prod_{\operatorname{Gal}(E/\Q_p)^{m-1}}\OO_E
\]
under the decomposition \eqref{eq:gl-tensor-decomposition}.
Thus $B_m$ is the integral closure of $A_m$.
The multiplication and literal unit $1$ here belong to the fixed
native field $E$. The normalization $I=p^{-1}\log\OO_E^\times$
specifies a lattice in that field; it does not replace its
multiplication by a structure transported through $\rho$.
In the declared additive tensor carrier put
\begin{equation}\label{eq:utp-point-source}
 Z_{j,\zeta_j}=1\otimes\cdots\otimes x_j,\qquad
 H_{j,\zeta_j}=
 \overline{\Span_{B_m}
   \{F^{\otimes m}(Z_{j,\zeta_j}):F\in\Gamma_E\}} .
\end{equation}
The literal ones agree with the single-slot embedding of
\cite[Proposition~3.1(ii), p.~93]{IUT3}. Choosing this
embedding does not itself prove global source membership.

\begin{theorem}[An attained point-source hull]\label{thm:utp-point-hull}
For every label and every finite torsion representative,
\begin{equation}\label{eq:utp-point-hull}
 H_{j,\zeta_j}=P_j
       =\beta^{\,e k_j-(e-1)m_j}B_{m_j}.
\end{equation}
Here the uniformizer exponent is taken in every field component.
For each tuple of representatives, one $F$ generates these
ideals simultaneously for all labels. With $\mu(B_m)=1$ and
$D_m=d^m$,
\begin{equation}\label{eq:utp-point-volume}
 D_{m_j}^{-1}\log\mu(H_{j,\zeta_j})
                       =(m_j\kappa-k_j)\log p .
\end{equation}
\end{theorem}
\begin{proof}
Every $F$ preserves $I$ and $p^{k_j}I$. Since
$x_j\in p^{k_j}I$ and $1\in I$, each tensor component
has valuation at least $k_j-m_j\kappa$. The component
embeddings preserve valuation, giving containment in $P_j$.
The common $F$ from Theorem~\ref{thm:utp-common-arrow} gives
$v_p(Fx_j)=k_j-\kappa$ and $v_p(F1)=-\kappa$.
Its tensor has the required valuation in every component,
so it is a generator of $P_j$ times an element of
$B_{m_j}^{\times}$. This proves the reverse inclusion
before closure; the fractional ideal is closed.
It is also the direct specialization $u=1$, $t=T_j$ of
Proposition~\ref{prop:gl-block-hull}.

There are $d^{m-1}$ field components, each with residue degree
two. Hence $\log\mu(\beta^\nu B_m)
=-2d^{m-1}\nu\log p$. Division by $d^m$ proves
\eqref{eq:utp-point-volume}. No further division by $m_j$
is made.
\end{proof}

The entire finite orbit $\mu_{2\ell}r_0^{j^2}$ gives the
same hull. The source in this proof contains no arbitrarily
chosen $\OO_E^\times$-multiple and need not contain
$r_0^{j^2}B_m$ or a product of fractional ideals.

\begin{proposition}[A comparison for each arrow]
\label{prop:utp-logarithmic-bridge}
For $0\ne x\in\OO_E$, set
$\lambda(x)=p^{-1}\log(1+px)$, $r=v_p(x)$ and
$k=\lfloor r+\kappa\rfloor$. For every $F\in\Gamma_E$,
\begin{equation}\label{eq:utp-tail}
 \begin{gathered}
 v_p(\lambda(x)-x)\ge1+2r\ge k+1-\kappa,\qquad
 \lambda(x)-x\in p^{k+1}I,\\
                       v_p(F(\lambda(x)))=v_p(F(x)).
 \end{gathered}
\end{equation}
Thus, for $t_j=\lambda(x_j)$ and $u=\lambda(1)$, the
point tensors with factors $(x_j,1,\ldots,1)$ and
$(t_j,u,\ldots,u)$ generate the same principal
$B_{m_j}$-ideal after each individual $F^{\otimes m_j}$.
Both orbit hulls are $P_j$, with the same common attainer.
\end{proposition}
\begin{proof}
For $n\ge2$, the $n$-th term of $\lambda(x)-x$ has valuation
$(n-1)+nr-v_p(n)$. For odd $p$, $v_p(n)\le n-2$,
so this is at least $1+2r$. Convergence and
ultrametricity give the first inequality; the second
uses $k\le r+\kappa$ and $r\ge0$.
By the definition of $k$, $x\in p^kI\setminus p^{k+1}I$.
Since $F$ is an automorphism of this filtered lattice,
\[
 v_p(Fx)<k+1-\kappa
                    \le v_p(F(\lambda(x)-x)).
 \]
Strict ultrametricity proves the valuation equality.
Also $u\in\Z_p^\times$, so $\Q_p$-linearity gives
$F(u)=uF(1)$. All component valuations of the two
tensors therefore agree, proving the ideal assertions.
Finally $(t_j-x_j)/p^{k_j}\in pI$, so the common
attainer for $T_j$ also works for $t_j/p^{k_j}$.
\end{proof}

Here $\lambda(x)$ is the coefficient of the separately
defined class $\operatorname{Kum}_p(1+px)$ in
$\rho=p^{-1}\log_{\rm BK}^{\rm std}$. It is not the
theta evaluation class. In particular, exact zero trace
was proved for $T_j$, not for $t_j$.
If all $m_j$ coordinates of this unit-cohomology packet
are changed from $\rho$ to $\log_{\rm BK}^{\rm std}$,
the tensor is multiplied by $p^{m_j}$. Consequently
\begin{equation}\label{eq:utp-standard-scale}
 \begin{gathered}
 P_j^{\rm std}=p^{m_j}P_j
                 =\beta^{\,e k_j+m_j}B_{m_j},\\
 D_{m_j}^{-1}\log\mu(P_j^{\rm std})
                       =-(k_j+m_j/e)\log p .
 \end{gathered}
\end{equation}
The change is $-m_j\log p$. This rescales every coordinate
of the specified cohomological source; it does not give the
original theta function an additional factor $p^{m_j}$.

\subsection{Synchronized source conditions}

The constructions distinguish
$\operatorname{Kum}(x_j)$, the numerical scalar $x_j$,
and $\operatorname{Kum}_p(1+px_j)$.
Their valuations already separate the first and third:
$v_p(x_j)=2j^2/\ell>0$, whereas $v_p(1+px_j)=0$.
Under normalized logarithm the scalar $x_j$ instead comes
from the unit $\exp(px_j)$. Equal additive hulls therefore
do not identify the multiplicative inputs.

The LGP construction uses the prescribed log-link pullback
before the single-slot insertion
\cite[Propositions~3.4(ii) and 3.5, pp.~102--106]{IUT3}.
A tautological local logarithmic marking can give this
pullback the numerical coordinate $x_j$. It neither makes
the log-link a ring homomorphism from the predecessor's
original field nor permits independent choices at all labels.
The theater log-link comes from one isomorphism and retains
its synchronization
\cite[Proposition~1.3(i) and Remark~1.3.1, pp.~42--43]{IUT3}.

Theorem~\ref{thm:utp-point-hull} alone does not establish
membership in the possible-image family of
\cite[Corollary~3.12, pp.~173--175]{IUT3}.
Section~\ref{sec:uniform-identity-membership} supplies one
globally synchronized copy representative and proves the
canonical local inclusion in its fixed base branch.
Compatibility of the complete family of log-Kummer maps
and reference data with the unindetermined $q$-pilot,
including the required Ind3 comparison, remains separate.
The full local Galois constructions used here are mathematical
inputs from Section~\ref{sec:full-galois-lifts}; their local
class field theory and source identifications are not newly
formalized in Lean by this section. No global abc conclusion
is asserted.
